%% SCRAP: hardware/platform-integration/l_4_re_blkio_endpoints
%% SOURCE: docs/working/hardware/platform-integration/l_4_re_blkio_endpoints.adoc
%% STATUS: WORKING
%% FITS: dev-guide/ch-l4re
%% EDITORIAL: lifted — prose rewritten to press voice

\section{Block I/O Porting Endpoints for L4Re}

The \texttt{blkio} subsystem maps cleanly onto L4Re because its northbound API
never changes. Porting to the microkernel means supplying new backends ---
dataspace-based or IPC-based --- behind the same six-function interface; the VM
and Forth block words keep calling exactly what they call today.

\subsection{The Stable Northbound API}

Every caller goes through this fixed surface, which constitutes the subsystem
ABI:

\begin{lstlisting}[language=C]
int  blkio_open (blkio_dev_t *dev, const blkio_vtable_t *vt, const blkio_params_t *p);
int  blkio_close(blkio_dev_t *dev);
int  blkio_read (blkio_dev_t *dev, uint32_t fblock, void *dst);
int  blkio_write(blkio_dev_t *dev, uint32_t fblock, const void *src);
int  blkio_flush(blkio_dev_t *dev);
int  blkio_info (blkio_dev_t *dev, blkio_info_t *out);
\end{lstlisting}

\subsection{In-Process Dataspace Backend}

When the VM can attach a dataspace directly, this is the fastest path --- no IPC
involved. The backend exposes \texttt{blkio\_l4ds\_state\_size()},
\texttt{blkio\_l4ds\_init\_state(...)}, and \texttt{blkio\_l4ds\_vtable()}, and
holds a small state record carrying the dataspace capability, the attached
virtual address, the mapped length, the block size, the block count, and a
read-only flag. The operations are direct: \texttt{open} attaches the dataspace
with \texttt{l4re\_rm\_attach()} and derives the block count;
\texttt{read}/\texttt{write} are \texttt{memcpy} against the mapped region;
\texttt{flush} is a no-op (or a dataspace operation if required); and
\texttt{close} detaches with \texttt{l4re\_rm\_detach()}.

\subsection{Out-of-Process Block Server Backend}

When storage is owned by another task --- an AHCI, NVMe, or eMMC server --- the
client-side backend exposes \texttt{blkio\_l4svc\_state\_size()},
\texttt{blkio\_l4svc\_init\_state(...)}, and \texttt{blkio\_l4svc\_vtable()},
talking to the server over a small opcode protocol:

\begin{lstlisting}[language=C]
enum bsvc_op {
  BSVC_INFO  = 1,
  BSVC_READ  = 2,
  BSVC_WRITE = 3,
  BSVC_FLUSH = 4,
  BSVC_PING  = 5
};
\end{lstlisting}

Two transports are available: a shared dataspace mapped by both server and
client (preferred), or inline IPC carrying the block payload in the message
(simpler but less efficient). The operations map onto the protocol directly ---
\texttt{read} sends the opcode and block number and the server fills the shared
buffer, \texttt{write} copies into the shared buffer before sending, and
\texttt{info} queries block size, count, and flags.

\subsection{Factory Helpers and CLI Semantics}

Backend selection can be hidden behind factory overloads ---
\texttt{blkio\_factory\_open\_l4ds(...)},
\texttt{blkio\_factory\_open\_file(...)}, and
\texttt{blkio\_factory\_open\_l4svc(...)} --- so \texttt{main.c} chooses a
backend from its arguments and environment. On L4Re the CLI follows the same
logic: \texttt{--disk-img=<path>} resolves either to a file server returning a
dataspace (the l4ds backend) or to a block server (the l4svc backend);
\texttt{--ram-disk=<MB>} allocates a RAM dataspace and uses the l4ds backend;
and with no disk image the system falls back to a RAM disk from the allocator.

\subsection{Composite Backends}

The intended layout dispatches by block number: blocks 0--1023 go to a RAM
backend and blocks 1024 and above go to a disk backend, both wrapped in a
composite vtable that routes on \texttt{fblock}. A later enhancement adds an
outer wrapper that packs three 1\,KiB blocks plus 1\,KiB of metadata into 4\,KiB
pages for device alignment --- still presenting the same five operations.

\subsection{Summary}

The L4Re port introduces two new backends (the in-process dataspace backend and
the IPC client/server backend) and, optionally, two factory helpers. Nothing
else moves: the VM continues to call \texttt{blkio\_read},
\texttt{blkio\_write}, and the rest. Backends are swapped, not call sites.
